Put the large centre at $(0,0)$ and $A=(-2,0)$. Let the third radius be $r$; the second centre is $(-2+3r,0)$. By reflection assume the third centre is $(u,r)$. Internal tangency to the large circle gives $u^2=4-4r$. External tangency to the second gives $(u+2-3r)^2=15r^2$. Eliminating $u$ yields $r^2(9r^2+84r-44)=0$. The positive root is $r=\dfrac{-14+4\sqrt{15}}{3}$. It satisfies the original distance equations with $u=\dfrac{3r^2+8r-4}{2-3r}$ and $0<3r<2$, so squaring has introduced no problem. The second radius is $3r=\sqrt{240}-14$, giving $\boxed{254}$.
\begin{asy}
import graph;import geometry;usepackage("amsmath");size(7cm);
real r=(-14+4sqrt(15))/3;pair O=(0,0),A=(-2,0),B=(2,0),D=(-2+3r,0),E=(sqrt(4-4r),r);draw(circle(O,2));draw(circle(D,3r));draw(circle(E,r));draw(A--B);draw(D--E,dashed);label("$A$",A,(-1,0));label("$B$",B,(1,0));label("$O$",O,(0,-1));label("$3r$",D,(0,1));label("$r$",E,(0,1));
\end{asy}

